\section{A trace family with potential automorphic support}
\label{sec:family}

We first place the given representation in a sufficiently large
fixed-determinant family.  We then put that entire family on completed
Hecke support after a totally real base change.  The representation
itself need not have an algebraic weight.  The distinction between its
residual representation and an auxiliary representation supplying the
determinant will be essential in the support argument.  For a matrix
representation, \emph{not virtually solvable} means that the identity
component of its algebraic monodromy group is not solvable.  This
property is unchanged on restriction to a finite index subgroup.

\subsection{Normalizing the determinant}

We will use the following twist property, proved in
Proposition~\ref{prop:hecke-twists}: at an odd tame level $N\geq5$, twisting
a Hecke coefficient point by any continuous character
$\eta:\Z_2^\times\to\OO'^\times$ gives another point of the same
Hecke algebra.  At a good prime $\ell$, this multiplies the trace by
$\eta(\eps(\operatorname{Frob}_\ell))$ and the determinant by its
square.  Thus it suffices to prove the theorem after any such twist
of $r$, and to undo the twist at the end.

\begin{lemma}\label{lem:normalization}
After a finite coefficient extension and a twist through $\eps$, one
may arrange
\[
 \det r=\chi=\delta\eps^w,
\]
where $\delta$ has finite order and is unramified at $2$.  The integer
$w$ may be chosen arbitrarily large in one parity class.  It may in
particular be chosen so that $\chi$ is the determinant of a regular
cuspidal modular representation that is not virtually solvable.
Oddness is preserved.  The character $\psi=\chi\eps$, viewed on
$\Z_2^\times$ by local reciprocity, is $z\mapsto z^{w+1}$.
\end{lemma}

\begin{proof}
Global class field theory separates the dyadic part of $\det r$ from
its odd conductor part:
\[
 \det r=\delta\,\alpha\circ\eps,
 \qquad \alpha:\Z_2^\times\longrightarrow\OO^\times.
\]
Here $\delta$ has finite order and odd conductor.  Indeed, at an odd
prime the pro-$\ell$ group of principal units has finite image in
$\OO^\times$, whose open subgroup is pro-$2$; only finitely many
primes occur.  The remaining unramified global character is trivial
by the class field theory of $\Q$.

Choose $w$ with $(-1)^w=\alpha(-1)$.  Write
$\Z_2^\times=\{\pm1\}\times5^{\Z_2}$.  The value $\alpha(5)$
reduces to one, since a pro-$2$ group has trivial image in
$\kk^\times$.  After finite coefficient extension choose
\[
 u^2=5^w\alpha(5)^{-1}.
\]
The element $u$ is integral and reduces to one.  Its powers $u^{2^n}$
tend to one, so $a\mapsto u^a$ defines a continuous character of
$\Z_2$.  Define $\beta(-1)=1$ and $\beta(5)=u$.  Then
$\alpha(z)\beta(z)^2=z^w$ for every $z\in\Z_2^\times$.
Replacing $r$ by $r\otimes(\beta\circ\eps)$ gives the asserted
determinant.  Since $\beta(-1)=1$, complex conjugation and oddness
are unchanged.  Local reciprocity identifies $\eps$ on dyadic units
with the identity character; $\delta$ is unramified there.

For completeness, the auxiliary modular representation can be chosen
without imposing the residual representation of $r$.  Oddness gives
$\delta(-1)=(-1)^{w+1}$.  Fix an odd level divisible by the conductor
of $\delta$, and enlarge it to a fine level.  The dimension formula
\cite[Section~III, Th\'eor\`eme~1]{CohenOesterle1977}
for cusp forms with this character grows linearly with the weight
$w+1$ along the indicated parity class.  The CM contribution at this
fixed level is bounded independently of the weight.  To see the
latter assertion, a CM eigenform comes from a Hecke character of an
imaginary quadratic field $K$; its level is
$|\operatorname{disc}K|\operatorname{Norm}(\mathfrak f)$, where
$\mathfrak f$ is the character conductor
\cite[Theorem~(3.4), Corollary~(3.5), and Theorem~(4.5)]{Ribet1977}.
Only finitely many pairs
$(K,\mathfrak f)$ occur at the fixed level.  For a fixed infinity
type the characters with conductor dividing a fixed ideal form, when
nonempty, a torsor under a finite ray class character group.  Oldform
multiplicities are bounded by the fixed level as well.  Thus a
sufficiently large $w$ gives a non-CM cuspidal eigenform of determinant
$\delta\eps^w$.  Its representation is regular and not virtually
solvable.  Indeed its restriction to every finite index subgroup is
semisimple; if its connected algebraic monodromy were solvable, its
unipotent radical would act trivially, so that connected group would
be a torus.  An open subgroup would then have abelian image, which
Ribet's CM criterion excludes
\cite[Propositions~(4.2) and~(4.4)]{Ribet1977}.
\end{proof}

We henceforth use the normalized $r$ and enlarge the finite
coefficient ring $\OO$ whenever necessary.  Every enlargement is
finite.  When a domain splits after coefficient extension, choose a
component through the specified coefficient point.  Finite flat
scalar extension makes its minimal prime contract to the old minimal
prime, so its dimension is unchanged.  Reduced support containments
also persist under this operation.  Thus the original coefficient
point and the dimension bounds below are retained.

\subsection{The dimension of the trace family}

\begin{lemma}\label{lem:large-family}
Fix a finite ramification set $S$ containing $2$, infinity, and the
ramified primes of the normalized $r$.  There is a reduced
horizontal framed deformation component $B$ of determinant $\chi$
through $r$.  Its completed trace image $D$ is a complete local
$\OO$-domain, contains the coefficient point of $r$, and satisfies
\[
 \dim B\ge6,\qquad \dim D\ge3.
\]
Moreover, $B$ has a coefficient point whose representation is not
virtually solvable.  Consequently the restriction of this family to
any finite extension of $\Q$ is generically absolutely irreducible.
\end{lemma}

\begin{proof}
Choose a stable lattice of $r$ and let $R^{\square,\chi}_{\bar r,S}$
be its unrestricted framed deformation ring with determinant $\chi$.
The residue here is the actual lattice reduction; no semisimplicity
or irreducibility is imposed on it.  These rings and their maps to
fixed-determinant pseudodeformation rings exist under Mazur's
finiteness condition, which holds for $G_{\Q,S}$
\cite[Proposition~3.1]{FM2}.

Complete at the characteristic-zero point of $r$.  Its real local
framed deformation ring is smooth of dimension two: an odd
involution has eigenvalues $1,-1$, and its conjugacy orbit has
stabilizer the two-dimensional diagonal torus.  For the global ring
over this real local ring, the relative relation estimate of
\cite[Proposition~3.6]{FM2}, with no imposed dyadic condition, is
\[
 g-r_{\mathrm{rel}}\ge
 3-h^0\bigl((\ad^0r)^\vee(1)\bigr)=3.
\]
The equality follows from absolute irreducibility.  In characteristic
zero the trace pairing identifies $\ad^0r$ with its dual; a nonzero
invariant in its cyclotomic twist would give a nonzero map
$r\to r\otimes\eps$.  Such a map is an isomorphism, whereas its
determinants would require $\eps^2=1$.  Every component of the
completed local ring therefore has dimension at least $2+3=5$.
The dimension comparison at a coefficient point in
\cite[Lemma~3.4]{FM2} adds one on returning to a horizontal integral
component.  Choose such a component through $r$ and give it the
reduced structure; its ring $B$ has dimension at least six.

Let $R^{\mathrm{ps},\chi}_{S}$ be the global pseudodeformation ring
of the semisimplified lattice reduction, and set
\[
 D=R^{\mathrm{ps},\chi}_{S}/
       \ker\bigl(R^{\mathrm{ps},\chi}_{S}\longrightarrow B\bigr).
\]
Write $T(g)\in D$ for its universal trace at $g$.  The ring $D$
is a horizontal complete local domain.  Its image in $B$ is
closed: complete local rings with finite residue field are compact,
and their continuous images in Hausdorff rings are closed.  The
generic representation on $B$ is absolutely irreducible, since the
irreducible locus contains $r$.  Forgetting its single frame loses
at most three dimensions, by \cite[Lemma~3.2]{FM2}.  Hence
$\dim D\ge\dim B-3\ge3$.

We justify the last assertion by bounding the trace loci of the
virtually solvable representations.  The reducible locus is closed
and proper on $D$, because it does not contain $r$.  An absolutely
irreducible virtually solvable two-dimensional representation is
induced from a character of a quadratic field, or has finite
projective image.  This follows by considering the connected
solvable algebraic subgroup: a noncentral torus has two eigenlines
permuted by the whole group, while a central connected subgroup gives
finite projective image.  A nontrivial connected unipotent subgroup
would have a unique invariant line and force reducibility.

Only finitely many quadratic fields $K$ can occur.  The corresponding
quadratic self-twist is unramified outside $S$, so this is
Hermite--Minkowski with bounded local degree.  For each $K$, the
maximal abelian pro-$2$ quotient of $G_{K,S}$ has rank at most two:
class field theory bounds its free part by the dyadic unit groups,
of total rank $[K:\Q]=2$; class groups and non-dyadic units contribute
only finite groups.  The invariant part under the quadratic
involution has rank at least one, from the cyclotomic extension.
If the inducing character is $\theta$, fixing the determinant fixes
$\theta\theta^s$ on $G_K$.  Thus at most one free character parameter
remains.  More explicitly, fixing the character on the image of
$1+s$ fixes a sublattice of rank at least one; the finite index in
its saturation contributes an integral, finite extension, not an
extra parameter.  The resulting complete character rings have
integral dimension at most two.  There are only finitely many
possible residual characters, since their unordered pair is
prescribed by $\bar r|_{G_K}$.

Their images in the trace space are closed of dimension at most two.
Indeed, each generating character value satisfies
\[
 X^2-T(h)X+\chi(h)=0\qquad(h\in G_K).
\]
Finitely many generators of the abelian character group, followed by
complete Nakayama, make each character ring finite over its trace
image.  This is the finite-character argument of
\cite[Proposition~3.1(3)]{FM2}.

For finite projective image, the classification of finite subgroups
of $\operatorname{PGL}_2$ leaves cyclic and dihedral groups, already
covered by reducibility and induction, and the three exceptional
groups $A_4,S_4,A_5$.  Their exponents divide $M=60$.
For any matrix with projective order dividing $M$, the ratio $\zeta$
of its eigenvalues satisfies $\zeta^M=1$, and
\[
 \frac{T(g)^2}{\chi(g)}=2+\zeta+\zeta^{-1}.
\]
Consequently $T(g)$ is a root of the monic polynomial
\[
 \prod_{\zeta^M=1}
  \bigl(X^2-\chi(g)(2+\zeta+\zeta^{-1})\bigr)\in\OO[X].
\]
Its coefficients belong to $\OO$ because the product over roots of
unity is Galois invariant and integral.  The fixed-determinant pseudo ring is
topologically generated by finitely many traces: traces generate it
by the determinant-of-a-sum identity, and finitely many suffice by
Nakayama on its cotangent space.  Imposing these monic equations on
such generators gives a finite $\OO$-algebra.  Thus the exceptional
finite-projective trace locus has integral dimension at most one.

The union of these finitely many closed loci cannot contain the
generic point of $D$, whose dimension is at least three.  Pull it
back to $B$ and use coefficient-point density
\cite[Lemma~3.3(2)]{FM2} to obtain a point outside it.  Restriction to
any finite index subgroup preserves the nonsolvable identity
component of its monodromy, so that point remains absolutely
irreducible.  The same is then true generically on the restricted
family and on every component containing that family.
\end{proof}

\subsection{Potential support with independently prescribed residue}

A closed trace locus over a totally real field is \emph{pro-modular}
if it lies in the spectrum of a residual factor of the completed
Hecke algebra of definite quaternionic forms, with the prescribed
central character and one tame level.  It is \emph{potentially
pro-modular} if this holds after a finite totally real solvable
extension, split completely above $2$.  All containments here use
the reduced closed trace images.  This is the support convention of
\cite[Section~2.3]{FM2}.

The potential pro-modularity theorem in \cite[Theorem~7.1]{FM2}
assumes a regular de Rham target.  Our family need not contain such
a target at this stage.  We therefore use its curve and connectedness
arguments with the determinant and residue specified separately.

\begin{proposition}\label{prop:potential-support}
Let $B$ be the family in Lemma~\ref{lem:large-family}, with the
normalized determinant of Lemma~\ref{lem:normalization}.  There is a
solvable totally real Galois extension $F/\Q$, of even degree and
completely split at $2$, such that the closed trace image of
$B|_{G_F}$ is pro-modular at one fixed tame level and central
character $\psi|_{G_F}$.  This assertion includes scalar and
reducible semisimple residual representations.
\end{proposition}

\begin{proof}
The determinant and the residual representation enter this argument
separately. Lemma~\ref{lem:normalization} supplies a determinant arising
from a regular modular representation; the modular seed below will
instead lift the residual representation of $B$.
We connect a characteristic-two curve on the seed's Hecke support
to one in an unrestricted locus containing the restricted family.
The constraints are used only to select and connect these curves.
Once the second curve has potential Hecke support,
localized propagation puts the entire unrestricted locus on that support.

We use the following precise propagation input.
The localized propagation theorem \cite[Theorem~5.1]{FM2} takes a
totally real solvable field of even degree, completely split at $2$,
a finite allowed set, a determinant of the form $\chi$ supplied by
a regular target, and, separately, a semisimple residual
representation of determinant $\bar\chi$.  If a characteristic-two
curve in this pseudodeformation space is potentially pro-modular,
has non-virtually-solvable generic representation, and has finite
local images at the allowed non-dyadic places, then every
irreducible closed locus containing that curve is potentially
pro-modular.

Choose a regular cuspidal modular lift of this semisimple residue,
using \cite[Lemma~7.2]{FM2}.  That lemma includes character sums and
scalar residue.  Enlarge the initial allowed set to contain its
tame primes and those of $B$, $\chi$, and the auxiliary regular
representation, before making any field choices.  Field preparation \cite[Proposition~6.1]{FM2} provides
abelian totally real $2$-extensions
\[
 \Q\subset A\subset F_0
\]
in which $2$ splits completely.  Put $d_A=[A:\Q]$ and
$d=[F_0:\Q]$.  The numbers $u_A,u_0$ of allowed non-dyadic primes
are bounded by a constant $U$ while both degrees grow.  Absolutely
irreducible residue remains absolutely irreducible.  In the
reducible case write $\bar r=\bar\chi_1\oplus\bar\chi_2$ and
$\alpha=\bar\chi_1/\bar\chi_2$.  Distinct characters remain
distinct, and restriction from $A$ to $F_0$ is injective on the
residual extension groups for $\alpha$ and $\alpha^{-1}$.

When the global residue is reducible and its dyadic characters coincide,
the connectedness argument will
keep track of chosen local invariant lines at selected dyadic places.
These places are the marks. Field preparation lets us make their
number large while keeping their proportion among all dyadic places
small.

There are no marks if the residue is absolutely irreducible or
$\alpha|_{G_{\Q_2}}\ne1$.  Otherwise field preparation specifies
one dyadic place $v_0$ of $A$, and we mark all its $m=[F_0:A]$
extensions to $F_0$.  Thus $m/d=1/d_A$; put $m=0$ in the unmarked
case.  We choose the degrees so that
\begin{equation}\label{eq:family-degrees}
 d>u_0+4,\qquad
 m>u_0+4\ \text{if marked},\qquad
 2d-2m-4-u_0>d.
\end{equation}
These are the degree conditions in \cite[Section~7.2]{FM2}.
Any further fixed lower bound on $d_A$ is compatible with them.

We now construct two large loci over $A$.  Solvable base change and
Jacquet--Langlands transfer the modular seed to definite forms.
Regularity keeps this base change cuspidal
\cite[Lemma~4.6]{FM2}.  Its central character and the prescribed
$\psi$ agree modulo $\varpi$.  At a sufficiently deep tame level,
the integral function description and exact reduction of completed
forms therefore give a nonzero residual summand with central
character precisely $\psi$; this is the seed construction in
\cite[Section~7.4]{FM2}.  A component $C_{\mathrm{seed}}$ of its
Hecke support has
\[
 \dim C_{\mathrm{seed}}\ge1+2d_A
\]
by \cite[Proposition~4.4]{FM2}.

Independently choose an unrestricted fixed-determinant framed
component containing the restricted family $B|_{G_A}$, by taking a
minimal prime below its defining prime.  Its generic representation
is absolutely irreducible by Lemma~\ref{lem:large-family}.
The relative presentation over the real local factors gives the
component dimension at least
$1+2d_A+3-h^0((\ad^0)^\vee(1))$; the invariant term is at most three.
Forgetting the frame loses at most three more dimensions.
Thus its trace image $C_{\mathrm{fam}}$, which contains the
restricted trace family, has
\[
 \dim C_{\mathrm{fam}}\ge2d_A-2
\]
by \cite[Proposition~3.6, Lemma~3.2, and Section~7.4]{FM2}.

In both loci impose $\varpi=0$ and the residual trace of one
Frobenius element at each allowed non-dyadic place.  In the marked
case impose also the constant residual local pseudorepresentation
at $v_0$.  The latter costs at most $\kappa$ equations, where
$\kappa$ is the number of generators of the closed-point ideal in
one fixed dyadic pseudo ring; it is independent of $d_A$.  Put
$\kappa=0$ when unmarked.  Taking
$d_A>U+\kappa+3$, the cut in $C_{\mathrm{fam}}$ has dimension at
least
\[
 2d_A-3-u_A-\kappa>d_A;
\]
the cut in $C_{\mathrm{seed}}$ has at least as large a lower bound.
Over the abelian field $A$, the characteristic-two reducible locus
has dimension at most one, the dihedral locus at most $d_A$, and
finite-projective-image curves are constant
\cite[Lemma~3.8]{FM2}.  Curve avoidance
\cite[Lemma~3.3(1)]{FM2} therefore gives a
non-virtually-solvable characteristic-two curve in each cut.
The Frobenius conditions make their allowed non-dyadic local images
finite by \cite[Lemma~3.9]{FM2}.  Denote these curves by
$C_s$ and $C_f$.

It remains to transfer support from $C_s$ to $C_f$.  Over $F_0$,
consider the characteristic-two framed representation scheme of
the fixed semisimple residue, adjoining an invariant projective
line at every mark.  A closed residual point is called
\emph{transverse} if none of these marked lines is globally
invariant; without marks this is an empty condition.
We check that both curves enter transverse residual charts.
Without marks, any stable lattice works.  With marks, the local
pseudo at $v_0$ is constant and scalar after a character twist.
The curve representation thus has a local invariant line, after a
finite extension of its field.  For a nonzero vector $v$ on that
line, over the normalized curve ring $J=\kk'[[t]]$ form
\[
 \mathcal V=\sum_{g\in G_A}J\rho(g)v.
\]
Compactness bounds this module in a lattice, and absolute
irreducibility makes it a full stable lattice.  The vector $v$ is
primitive: otherwise every translate would belong to
$t\mathcal V$, and Nakayama would give $\mathcal V=0$.
Its reduced orbit spans $\mathcal V/t\mathcal V$, so its reduction
belongs to no global invariant line.

This avoidance survives restriction to $F_0$.  A nonsplit extension
of distinct residual characters remains nonsplit by the restriction
injectivity, while a split representation retains its two distinct
character lines.  For globally scalar semisimplification, the orbit
condition rules out a split scalar reduction; its nonzero additive
extension remains nonzero by the same injectivity.  These statements
hold after any residual field extension.  Finally, $F_0/A$ is
Galois, so transporting the local line to all conjugate marks
preserves avoidance of every global invariant line.  Both curves
therefore pass through transverse residual points, including in
the scalar residual case.

We recall what the two connectedness inputs prove, to make the
support transfer explicit.  At a transverse point,
\cite[Lemma~7.3]{FM2} gives components of the completed chart of
dimension at least $2d-m$, with successive intersections in a
connecting chain of dimension at least $2d-m-1$.  The reducible
locus has strictly smaller dimension.  For an intersection
component of maximal dimension, forgetting its frame and marked
lines loses at most $3+m$ dimensions, and forgetting the characters
on those lines is finite.  Its trace image therefore has dimension
at least $2d-2m-4$.  The $u_0$ Frobenius trace conditions leave
dimension greater than $d$ by \eqref{eq:family-degrees}, so the same
bad-locus avoidance supplies a non-virtually-solvable curve with
finite bad local images.  Localized propagation carries support
across each such intersection.  This is
\cite[Lemma~7.4]{FM2}.

The transverse residual fiber is geometrically connected by
\cite[Lemma~7.5]{FM2}.  In particular this includes scalar residue:
after fixing the unique global invariant line, the nonzero additive
extension classes and the transverse marked lines give a connected
parameter space.  The finite graph of components meeting this fiber
therefore joins the seed chart to the family chart.  Starting from
$C_s$, which is pro-modular, the preceding chart propagation makes
$C_f|_{G_{F_0}}$ potentially pro-modular.  Hence $C_f$ over $A$ is
potentially pro-modular.  Apply localized propagation now over $A$
to the \emph{unrestricted} locus $C_{\mathrm{fam}}$ containing
$C_f$.  The entire locus, and therefore the prescribed family,
acquires Hecke support.  Marks have served only to connect curves
in characteristic two; they impose no condition on $B$.

Only finitely many components, charts, and propagation steps occur.
Taking the common Galois compositum of the resulting fields and a
common deeper tame level, as in \cite[Lemma~4.6 and Section~7.4]{FM2},
gives one extension $F/\Q$ with all the stated properties.
\end{proof}

\subsection{The diagonal local parameter and completed forms}

Write $D$ for the trace domain of Lemma~\ref{lem:large-family} and
choose $F$ from Proposition~\ref{prop:potential-support}.  Let
$\mathbb T$ be the residual completed Hecke algebra over $F$ just
obtained.  Support containment gives a continuous map
$\mathbb T\to D$: its image is the closed algebra of traces
restricted to $G_F$.  Indeed the Hecke kernel vanishes in this
reduced domain, by the definition of closed support.

Let $R$ be the local pseudodeformation ring of the residual dyadic
parameter, with determinant $\chi|_{G_{\Q_2}}\eps^2$.  This is
\emph{block normalization}: we twist each Galois parameter by
$\eps$, and the central character is
$\psi=\chi\eps$.  Because $2$ splits completely in $F$, each of
the $b=[F:\Q]$ dyadic restrictions of the global family has this
same parameter as a $D$-valued determinant.  The identifications use conjugate decomposition
groups; traces are unchanged by the conjugations.

For the corresponding compact dual block of
$G=\GL_2(\Q_2)$ with this central character, let $P$ be the direct
sum of one projective envelope of the dual of each simple object,
and put $\mathcal E=\End(P)$.  Coefficients are enlarged so that
the block is absolutely split.  The projective-generator functor
$\Hom(P,-)$ takes values in right $\mathcal E$-modules; its inverse
is completed tensor product with $P$.  We use the exact compact
block equivalence here, and recall its finer local properties in
Section~\ref{sec:diagonal}.

The next proposition puts the trace family and these local actions on
one coefficient module. Its finiteness over $D$ and $R$ will let us
compare specialization with admissibility.

\begin{proposition}\label{prop:family}
For the normalized representation $r$, the preceding constructions
give a horizontal complete local domain $D$ of dimension at least
three containing its coefficient point, and continuous maps
\[
 \mathbb T\longrightarrow D,\qquad R\longrightarrow D
\]
that make $D$ finite over both source rings.  There is a right
$\mathcal E^{\widehat\otimes b}$-module $L_0$, finite and faithful
over $\mathbb T$, whose inverse under the product block equivalence
is the residual summand of compact definite forms.  The module
\[
 L=\bigl(L_0\otimes_{\mathbb T}D\bigr)
       \big/\bigl(\text{$D$-torsion}\bigr)
\]
is nonzero, finite and torsion-free over $D$.  All $b$ central
$R$-actions on $L$ are the same map $R\to D$.  Its inverse under
the product block equivalence is finite over the completed group
algebra of a sufficiently small product of determinant-one compact
subgroups.
\end{proposition}

\begin{proof}
Only finiteness and the module assertions remain to be proved.
Let $D_F$ be the closed restricted trace image in $D$.  If
$n=[F:\Q]$, then $g^n\in G_F$ for every $g\in G_\Q$.
The trace-of-powers recurrence
\[
 P_0(X,a)=2,\quad P_1(X,a)=X,\quad
 P_j(X,a)=XP_{j-1}(X,a)-aP_{j-2}(X,a)
\]
gives a monic polynomial $P_n(X,a)$ of degree $n$, with
\[
 P_n(T(g),\chi(g))=T(g^n)\in D_F.
\]
Choose finitely many traces topologically generating $D$ over
$\OO$.  Each is integral over $D_F$, so the subalgebra they generate
is finite over $D_F$.  It is compact and hence closed in $D$;
since it also contains a dense subalgebra, it equals $D$.
Thus $D$ is finite over $D_F$, and therefore over $\mathbb T$.

The local-global block comparison \cite[Proposition~4.2]{FM2}
makes $\mathbb T$ finite over
$R^{\widehat\otimes b}$, with its actual Galois restriction action
in block normalization.  On $D$ this action factors through
\[
 R^{\widehat\otimes b}\longrightarrow R,
 \qquad a_1\otimes\cdots\otimes a_b\longmapsto a_1\cdots a_b,
\]
because all dyadic parameters agree.  Transitivity of finiteness
then proves that $D$ is finite over the single ring $R$.

Let $M$ be the compact definite-forms summand.  The same proposition
gives
\[
 L_0=\Hom_{G^b,\mathrm{cts}}(P^{\widehat\otimes b},M),
\]
finite and faithful over $\mathbb T$, with the asserted compatible
product action.  The compact module $M$ is admissible, equivalently
finite over a sufficiently deep determinant-one completed group
algebra, by \cite[Proposition~4.1]{FM2}.

At the generic point of $D$, the base change of $L_0$ is nonzero.
Indeed, if $\mathfrak q$ is its contraction to $\mathbb T$,
faithfulness and finiteness give
$\operatorname{Supp}_{\mathbb T}L_0=\Spec\mathbb T$; Nakayama
then gives $L_0\otimes_{\mathbb T}\kappa(\mathfrak q)\ne0$.
Extending this residue field to $\Frac(D)$ preserves nonvanishing.
Quotienting by $D$-torsion consequently leaves a nonzero finite
$D$-module $L$.  The torsion submodule is stable under every
$\mathcal E$-action because those actions commute with $D$.
Since $D$ is horizontal, $L$ is also $\OO$-torsion-free.

Finally, finiteness of $D$ over $\mathbb T$ makes
$L_0\otimes_{\mathbb T}D$, and then $L$, a quotient of a finite
direct sum of copies of $L_0$ as product-block modules.  These are
maps of compact modules with closed images.  Exactness of the
product block equivalence makes the inverse object of $L$ a
quotient of the corresponding finite sum of copies of $M$.
It is therefore admissible as asserted.  This supplies product
admissibility; passing to a single factor is the separate argument
of the next section.
\end{proof}
